\documentclass[12pt,a4paper]{article}
\usepackage[utf8]{inputenc}
\usepackage[catalan]{babel}
\usepackage{graphicx}
\usepackage{geometry}
\usepackage{fancyhdr}
\usepackage{tabularx}
\usepackage{lastpage}
\usepackage{enumitem}
\usepackage{amsmath}
\usepackage{tikz}

% 1. MARGES I CAPÇALERA
\geometry{top=4cm, bottom=2cm, left=2cm, right=2cm, headheight=2.5cm}

\pagestyle{fancy}
\fancyhf{}
% Assegura't que tens els fitxers d'imatge o comenta aquestes línies si no els tens
\fancyhead[L]{\includegraphics[height=1.8cm]{Logo_Departament.png}} 
\fancyhead[R]{
\includegraphics[height=1.8cm]{Logo_Institut.png}}

\fancyfoot[L]{}
\fancyfoot[C]{Pàgina \thepage\ de \pageref{LastPage}}
\fancyfoot[R]{}

\renewcommand{\headrulewidth}{0.4pt}

\begin{document}

% --- CAPÇALERA D'IDENTIFICACIÓ ---
\begingroup
\renewcommand{\arraystretch}{2.2}
\begin{tabularx}{\textwidth}{|X|p{4cm}|}
    \hline
    \textbf{Nom i cognoms:} & \textbf{Grup:} \\ \hline
    \textbf{Activitat:} Introducció a l'estadística & \textbf{Data:} \\ \hline
\end{tabularx}
\endgroup

\vspace{0.4cm}

% --- RECORDATORI ---
\noindent\fbox{
\begin{minipage}{\dimexpr\textwidth-2\fboxsep-2\fboxrule\relax}
\textbf{Recordatori:}
\begin{itemize}[noitemsep, topsep=0pt]
    \item \textbf{Mitjana ($\bar{x}$):} Suma de totes les dades dividida pel nombre total de dades ($N$).
    \item \textbf{Mediana ($Me$):} El valor central un cop ordenades les dades de menor a major.
    \item \textbf{Moda ($Mo$):} El valor o dada que més es repeteix.
    \item \textbf{Freqüència absoluta ($f_i$):} Nombre de vegades que apareix una dada.
    \item \textbf{Freqüència relativa ($h_i$ o $f_r$):} Quocient entre la freqüència absoluta i el total ($f_i / N$).
\end{itemize}
\end{minipage}}

\vspace{0.4cm}

% --- ACTIVITAT 1 ---
\noindent\textbf{Activitat 1:} Hem preguntat a 20 alumnes quantes hores han dormit aquesta nit. Les respostes han estat: \\
\textbf{6, 7, 8, 7, 6, 9, 7, 8, 6, 7, 8, 7, 9, 6, 7, 8, 7, 6, 8, 7}

\begin{enumerate}[label=\alph*)]
    \item \textbf{Completa la taula de freqüències:}
    
    \vspace{0.1cm}
    \renewcommand{\arraystretch}{1.5}
    \begin{table}[h]
    \begin{tabular}{|l|c|c|c|c|c|}
        \hline
        \textbf{Hores de son ($x_i$)} & 6 & 7 & 8 & 9 & \textbf{Total ($N$)} \\ \hline
        \textbf{Freq. absoluta ($f_i$)} & \hspace{1cm} & \hspace{1cm} & \hspace{1cm} & \hspace{1cm} & \\ \hline
        \textbf{Freq. relativa ($h_i$)} & & & & & \\ \hline
    \end{tabular}
    \end{table}

    \item \textbf{Calcula els paràmetres estadístics:}
    
    \vspace{0.2cm}
    $\bar{x} =$ \hspace{3cm} $Me =$ \hspace{3cm} $Mo =$
    
    \vspace{0.4cm}
    \item \textbf{Dibuixa el gràfic de barres:}
    
    \vspace{0.1cm}
    \begin{tikzpicture}[scale=0.6]
        % Eixos
        \draw[thick, ->] (0,0) -- (6,0) node[right] {$x_i$ (hores)};
        \draw[thick, ->] (0,0) -- (0,8.5) node[above] {$f_i$};
        
        % Quadricula suau
        \draw[very thin, gray!50, step=1] (0,0) grid (5.8,8.2);
        
        % Marques eix Y
        \foreach \y in {1,2,3,4,5,6,7,8} {
            \draw (-0.1,\y) -- (0.1,\y) node[left=0.2cm] {\y};
        }
        % Marques eix X
        \foreach \x/\label in {1/6, 2.5/7, 4/8, 5.5/9} {
            \draw (\x,-0.1) -- (\x,0.1) node[below=0.2cm] {\label};
        }
    \end{tikzpicture}
\end{enumerate}

\vspace{0.1cm}

\newpage
% --- ACTIVITAT 2 ---
\noindent\textbf{Activitat 2:} En una classe s'ha fet una enquesta ràpida sobre el nombre de llibres llegits el darrer mes. Les 20 respostes són:\\
\textbf{1, 2, 1, 3, 0, 2, 1, 4, 2, 1, 2, 0, 3, 1, 2, 2, 1, 0, 3, 2}

\vspace{0.2cm}
\begin{enumerate}[label=\alph*)]
    \item \textbf{Construeix la taula de freqüències:}
    
    \vspace{0.1cm}
    \renewcommand{\arraystretch}{1.5}
    \begin{table}[h]
    \begin{tabular}{|l|c|c|c|c|c|c|}
        \hline
        \textbf{Llibres ($x_i$)} & \hspace{0.8cm} & \hspace{0.8cm} & \hspace{0.8cm} & \hspace{0.8cm} & \hspace{0.8cm} & \textbf{Total ($N$)} \\ \hline
        \textbf{Freq. absoluta ($f_i$)} & & & & & & \\ \hline
        \textbf{Freq. relativa ($h_i$)} & & & & & & \\ \hline
    \end{tabular}
    \end{table}
    
    \item \textbf{Calcula la mitjana, la mediana i la moda:}
    
    \vspace{0.6cm}
    $\bar{x} =$ \hspace{3cm} $Me =$ \hspace{3cm} $Mo =$
    \vspace{0.6cm}
    
    \item \textbf{Dibuixa el corresponent gràfic de barres} (recorda posar nom als eixos):
    \vspace{4.5cm} % Espai en blanc per fer el gràfic a mà
\end{enumerate}

\end{document}